\begin{textsample}{POS Dim 2 – df – Score 17.00 – df/df\_em\_53.txt}  \label{ex:f2_pos_130}
CARACTERIZAÇÃO DA EDUCAÇÃO \textbf{PROFISSIONAL} E TECNOLÓGICA ( EPT ) Conforme disposto na LDB, com base nas alterações introduzidas pela Lei nº 11.741/2008, a Educação Profissional e Tecnológica ( EPT ) integra -se aos diferentes níveis e modalidades de educação e às dimensões do trabalho, da ciência, da tecnologia e da cultura, consistindo em uma modalidade da Educação Básica e Superior, em articulação com as modalidades de Educação de Jovens e Adultos ( EJA ), Educação Especial e Educação a Distância ( EaD ) ( BRASIL, 2008b ).
Enquanto modalidade da Educação Básica, a EPT \textbf{destina} -se à formação \textbf{profissional} dos jovens e dos adultos, com o objetivo de prepará -los para a vida, incluindo a sequência dos estudos no nível \textbf{técnico}, tecnológico e superior, bem como inseri -los no mundo do trabalho ( BRASIL, 2008b ).
Para fins de alinhamento conceitual, a Resolução nº 3/2018 do Conselho Nacional de Educação ( CNE ), a qual \textbf{atualiza} as \textbf{Diretrizes} Curriculares Nacionais para o Ensino Médio ( DCNEM ), estabelece o conjunto de termos e conceitos próprios relacionados ao \textbf{itinerário} de Formação \textbf{Técnica} e Profissional ( BRASIL, 2018b ), acrescido de terminologia definida pelo Conselho de Educação do Distrito Federal ( CEDF ), consoante o quadro resumido abaixo ( \textbf{DISTRITO} \textbf{FEDERAL}, 2019b ) : | TERMO | SIGNIFICADO | |---|---| | Ambientes Simulados | Ambientes pedagógicos que possibilitam o desenvolvimento de atividades práticas da aprendizagem \textbf{profissional}. | | Formações Experimentais | Formações autorizadas pelos respectivos sistemas de ensino, nos termos de sua \textbf{regulamentação} específica, que ainda não constam no CNCT. | | Aprendizagem Profissional | Formação técnico-profissional compatível com o desenvolvimento físico, moral, psicológico e social do jovem de 14 a 24 anos de idade. | | \textbf{Qualificação} \textbf{Profissional} | Processo ou resultado de formação e desenvolvimento de competências de um determinado \textbf{perfil} \textbf{profissional}, definido no mercado de trabalho. | | \textbf{Habilitação} \textbf{Profissional} \textbf{Técnica} de Nível Médio | \textbf{Qualificação} \textbf{profissional} formalmente reconhecida por meio de diploma de conclusão de \textbf{curso} \textbf{técnico}, com validade nacional. | | Programa de Aprendizagem | Arranjos e combinações de \textbf{cursos} que, articulados e com os devidos aproveitamentos curriculares, possibilitam um \textbf{itinerário} \textbf{formativo}. | | \textbf{Certificação} Intermediária | \textbf{Certificação} de \textbf{qualificação} para o trabalho quando a formação for estruturada e organizada em etapas com terminalidade. | | \textbf{Certificação} \textbf{Profissional} | Processo de avaliação, reconhecimento e \textbf{certificação} de saberes adquiridos na Educação Profissional, inclusive no trabalho. | | Eixo Tecnológico | Agrupamento sistematizado de conhecimentos, com interdependências entre as áreas científicas e culturais, comuns a grandes ramos de conhecimentos tecnológicos e áreas \textbf{profissionais}. | | Saída Intermediária | Etapa de um \textbf{Curso} \textbf{Técnico} de Nível Médio com caráter de terminalidade, correspondente a um \textbf{curso} de Formação Inicial e Continuada ( FIC ), ou de \textbf{Qualificação} \textbf{Profissional}, composta por uma ou mais unidades curriculares, definida(s) no Plano de \textbf{Curso}. | | Arranjos Produtivos Locais ( APL ) | Conjunto de agentes econômicos, políticos e sociais localizados no mesmo território, desenvolvendo atividades econômicas correlatas e que apresentam vínculos expressivos de produção, interação, cooperação e aprendizagem ( OLIVEIRA et al., 2017 ). | Do ponto de vista da \textbf{oferta}, a EPT compreende os \textbf{cursos} de nível médio e os \textbf{cursos} de \textbf{qualificação} \textbf{profissional} ou de formação inicial e continuada, conforme disposto abaixo :

% matched lemmas: atualizar, certificação, curso, destinar, diretriz, distrito, federal, formativo, habilitação, itinerário, oferta, perfil, profissional, qualificação, regulamentação, técnico
\end{textsample}
